\makeatletter \def\input@path{{../../}} \makeatother
\documentclass[11pt]{article}

\usepackage[utf8]{inputenc}
\usepackage{graphicx}
\usepackage[dvipsnames]{xcolor}
\usepackage{color}
\usepackage[margin=1in]{geometry}
\usepackage{hyperref}
\usepackage{tikz}
\usepackage{amsmath}
\usepackage{commath}
\usepackage{relsize}
\usepackage{centernot}
\usepackage{amssymb}
\usepackage{amsfonts}
\usepackage{graphicx}
\usepackage{textcomp}
\usepackage{gensymb}
\usepackage{enumitem}
\usepackage{siunitx}
\usepackage{fancyhdr}
\usepackage{floatrow}
\usepackage{wrapfig}
\usepackage[labelfont=bf]{caption}
\usepackage{mathtools}
\usepackage{pgfplots}
\usepackage{multicol}
\usepackage[misc]{ifsym}
\usepackage{tabularx}
\usepackage{empheq}
\usepackage{tkz-tab}
\usepackage{xpatch}
\usepackage{caption}
\usepackage{subcaption}
\usepackage{tabularray}
\usepackage{esint}
\usepackage{amsthm}
\usepackage{thmtools}
\RequirePackage[most]{tcolorbox}

\swapnumbers

\makeatletter
\declaretheoremstyle[
  headfont= \sffamily\bfseries\color{cyan!50!black},
  headpunct={\sffamily\bfseries.},
  postheadspace=0.5em,
  notefont=\sffamily\bfseries,
  headformat=\NAME\NUMBER\let\thmt@space\@empty\NOTE,
  bodyfont=\mdseries\slshape,
  spaceabove=12pt,spacebelow=12pt,
]{thmstyle}
\makeatother

\theoremstyle{thmstyle}
\declaretheorem[name=Definici\'on]{theorem}
\newtheorem{definition}[theorem]{Definici\'on}

\definecolor{lightcyan}{rgb}{0.88, 1.0, 1.0}
\colorlet{mythmback}{lightcyan!40!white}

\tcolorboxenvironment{theorem}{
  colback=mythmback,coltitle=blue,colframe=mythmback,
  left=0pt,right=0pt,
  top=0pt,bottom=0pt,
  before skip=10pt, after skip=10pt,
}

\xpatchcmd{\tkzTabLine}{$0$}{$\bullet$}{}{}
\tikzset{t style/.style={style=solid}}
\pgfplotsset{compat=newest}

\newtheorem{manualtheoreminner}{}
\newenvironment{manualtheorem}[1]{%
  \renewcommand\themanualtheoreminner{#1}%
  \manualtheoreminner
}{\endmanualtheoreminner}

\renewcommand{\pd}[2]{\frac{\partial #1}{ \partial #2}}
\newcommand{\td}[2]{\frac{\mathrm{d} #1}{ \mathrm{d} #2}}
\newcommand{\pdd}[2]{\frac{\partial^2 #1}{ \partial #2 ^2}}
\newcommand{\pddc}[3]{\frac{\partial^2 #1}{ \partial #2 \partial #3}}
\newcommand{\solution}[0]{\large{\textbf{Soluci\'on}}\normalsize}

\def\sca   #1{\mbox{\rm{#1}}{}}
\def\mat   #1{\mbox{\boldmath $\mathsf #1$}}
\def\vec   #1{\mbox{\boldmath $#1$}{}}
\def\ten   #1{\mbox{\boldmath $#1$}{}}
\def\ltr   #1{\mbox{\sf{#1}}}
\def\bltr  #1{\mbox{\sffamily{\bfseries{{#1}}}}}

\DeclareMathOperator{\dive}{div}
\DeclareMathOperator{\trace}{trace}
\DeclareMathOperator{\tr}{tr}
\DeclareMathOperator{\symm}{symm}
\DeclareMathOperator{\sk}{skew}
\DeclareMathOperator{\grad}{grad}
\DeclareMathOperator{\Grad}{Grad}
\DeclareMathOperator{\curl}{curl}
\DeclareMathOperator{\Curl}{Curl}
\DeclareMathOperator*{\argmin}{\arg\!\min}

\def\dx{\mbox{\(\,\mathrm{d}x\)}}
\def\ds{\mbox{\(\,\mathrm{d}s\)}}
\def\dS{\mbox{\(\,\mathrm{d}S\)}}
\def\dV{\mbox{\(\,\mathrm{d}V\)}}
\def\dv{\mbox{\(\,\mathrm{d}v\)}}
\def\R{\mathbb R}
\def\C{\mathbb C}
\def\N{\mathbb N}
\def\Q{\mathbb Q}
\def\Z{\mathbb Z}
\def\D{\mathcal D}

\usepackage{mdframed}
\mdfdefinestyle{MyFrame}{%
    linecolor=black,
    linewidth=1pt,
    outerlinewidth=1pt,
    roundcorner=20pt,
    innertopmargin=\baselineskip,
    innerbottommargin=\baselineskip,
    innerrightmargin=20pt,
    innerleftmargin=20pt,
    splittopskip=2\baselineskip,
    backgroundcolor=gray!50!white}

\setlength\textwidth{7in}
\setlength\parindent{0pt}
\oddsidemargin=-0.5cm
\pagestyle{fancy}
\newcommand*{\vertbar}{\rule[-1ex]{0.5pt}{2.5ex}}
\setlength{\headheight}{13.59999pt}

\renewcommand{\headrulewidth}{0.1pt}
\renewcommand{\footrulewidth}{0.1pt}
\renewcommand{\figurename}{Figura}
\renewcommand\refname{Referencias}
\renewcommand{\thesection}{\arabic{section}.}
\renewcommand{\thesubsection}{\thesection\arabic{subsection}.}
\renewcommand{\labelitemi}{{\raisebox{.3\height}{\scalebox{0.6}{$\blacklozenge$}}}}
\renewcommand*\contentsname{\'Indice}

\newcommand{\p}{\vspace{10pt}}

\AtBeginDocument{\vspace*{-3.2\baselineskip}}
\textheight=665pt

\lhead{\scshape Pontificia Universidad Cat\'olica de Chile}
\rhead{\scshape IMC UC}
\cfoot{\thepage}

\usepackage{footnote}
\usepackage{etoolbox}
\newtoggle{indefintion}
\togglefalse{indefintion}
\pretocmd{\footnote}{\iftoggle{indefintion}{\stepcounter{footnote}}{\relax}}{}{}

\begin{document}
    \begin{flushleft}
        \small
        \vspace{1pt}
        \textbf{IMT3130} \hfill
        \textbf{24/04/2026}
    \end{flushleft}
    \begin{center}
        \LARGE\textbf{Ayudant\'ia 6}\\
        \vspace{3pt}
        \normalsize\textbf{Elementos finitos $P^1$, lifting y ensamblaje}\\
        \normalsize Ayudante: Basti\'an Herrera \Letter{\hspace{1pt}: \texttt{biherrera@uc.cl}}
        \rule{\textwidth}{0.05mm}
    \end{center}

\section*{Problemas}

\begin{enumerate}[label=\textbf{\arabic*.}, leftmargin=*]

    \item \textbf{Elementos $P^1$ en 1D, matriz local y lifting discreto.} Sea $\Omega=(0,1)$ y considere una partici\'on uniforme $0=x_0<x_1<\cdots<x_N=1$, con $h=1/N$. Sea $P^1_h$ el espacio de funciones continuas y lineales por partes sobre la malla, y considere el problema
    \[
    \begin{aligned}
        -\kappa u''+\sigma u &= f &&\text{en }(0,1),\\
        u(0)&=g_0,\\
        u(1)&=g_1,
    \end{aligned}
    \]
    donde $\kappa>0$, $\sigma\ge 0$ y $f\in L^2(0,1)$. Defina
    \[
        V_h^0=\{v_h\in P^1_h:v_h(0)=v_h(1)=0\},
        \qquad
        V_h^g=\{v_h\in P^1_h:v_h(0)=g_0,\ v_h(1)=g_1\}.
    \]
    \begin{enumerate}[label=(\alph*)]
        \item Tome un lifting discreto $G_h\in P^1_h$ tal que $G_h(0)=g_0$ y $G_h(1)=g_1$, escriba $u_h=u_{0,h}+G_h$ con $u_{0,h}\in V_h^0$, y derive el sistema lineal satisfecho por los grados de libertad interiores de $u_{0,h}$.
        \item Para un elemento $K_e=[x_e,x_{e+1}]$, calcule la matriz elemental asociada a
        \[
            a(w,v)=\int_0^1 \kappa w'v'\,dx+\int_0^1 \sigma wv\,dx.
        \]
        \item Ensamble la forma tridiagonal de la matriz global reducida y explique c\'omo cambian la primera y la \'ultima ecuaci\'on interior cuando $g_0$ y $g_1$ no son nulos.
    \end{enumerate}

    \item \textbf{Ensamblaje 2D con condiciones Dirichlet y Robin.} Considere el cuadrado $\Omega=(0,1)^2$ con nodos globales
    \[
        x_1=(0,0),\qquad x_2=(1,0),\qquad x_3=(1,1),\qquad x_4=(0,1),
    \]
    dividido en los tri\'angulos $K_1=(x_1,x_2,x_3)$ y $K_2=(x_1,x_3,x_4)$. Considere el problema
    \[
    \begin{aligned}
        -\Delta u &= 1 &&\text{en }\Omega,\\
        \gamma_0u &=0 &&\text{en }\Gamma_D:=\{0\}\times(0,1),\\
        \partial_nu+\beta\gamma_0u &= r &&\text{en }\Gamma_R:=\{1\}\times(0,1),\\
        \partial_nu&=0 &&\text{en }\partial\Omega\setminus\overline{(\Gamma_D\cup\Gamma_R)},
    \end{aligned}
    \]
    donde $\beta\ge 0$ y $r\in\R$ son constantes. Use elementos finitos $P^1$ conformes y tome como nodos Dirichlet a $x_1$ y $x_4$.
    \begin{enumerate}[label=(\alph*)]
        \item Escriba la formulaci\'on de Galerkin discreta y construya los arreglos IEN, ID y LM para esta malla.
        \item Calcule las matrices elementales de volumen para el t\'ermino de Poisson, los vectores locales de carga asociados a $f\equiv 1$, y ensamble la matriz y el vector globales reducidos antes de incorporar Robin.
        \item Calcule la contribuci\'on de Robin sobre la arista $E_R=[x_2,x_3]$ y escriba el sistema lineal final. Explique brevemente qu\'e cambiar\'ia si los valores Dirichlet en $x_1$ y $x_4$ fueran no homog\'eneos.
    \end{enumerate}

\end{enumerate}

\end{document}
